\def\examid{26101200001}
\begin{document}
\begin{enumerate}
\item \multipart{1}{8}
The data set \code{fertil1} contains data from women in the US between the ages of 35 and 54. You can load the data into R using the \code{wooldridge} library using:

\medskip
\code{library(wooldridge)}\\
\code{dat = wooldridge::fertil1}
\medskip

The main outcome of interest is the variable \textit{kids}, which is the number of children the women has given birth to.\\
What is the mean of the variable \textit{kids} in the dataset?
\begin{opts}
\item \ca{2.7431}
\item 1.8002
\item 3.8678
\item 2.5798
\item None of the above
\end{opts}

\item How many missing values or NA's are there for the variable \textit{age} in the dataset?
\begin{opts}
\item \ca{0}
\item 265
\item 20
\item 33
\item None of the above
\end{opts}

\item Consider the following linear regression model.
\[\begin{aligned}
kids = \beta_0 + \beta_1(educ) + \beta_2(age) + \beta_3(age^2) + \beta_4(year) &+ \\
\beta_5(east) + \beta_6(northcen) + \beta_7(west) &+ \varepsilon
\end{aligned}\]
The model includes controls for education, age and age squared. The model also controls for \textit{year}, which is year. Finally, (\textit{east, northcen, west}), indicator variables for regional areas are included in the model.\\
Estimate the model above with OLS. What is the estimated coefficient on education?
\begin{opts}
\item $-0.06932$
\item $-0.13712$
\item $-0.15261$
\item $-0.17425$
\item \ca{$-0.12260$}
\end{opts}

\item What is the sum of squared errors of the model?
\begin{opts}
\item 6632.205
\item \ca{2809.876}
\item 5223.497
\item 861.390
\item None of the above
\end{opts}

\item Suppose we wanted to create a confidence interval for one of the parameters of the model. If we wanted to create a 97.50\% confidence interval, what critical value would we use?
\begin{opts}
\item 2.4559
\item 3.1180
\item \ca{2.2444}
\item 2.2664
\item None of the above
\end{opts}

\item Consider the null hypothesis that $H_0:\beta_{northcen}=0$. What is the strongest conclusion we can draw about this null hypothesis?
\begin{opts}
\item Reject the null at the 0.1\% level
\item Reject the null at the 1\% level
\item \ca{Reject the null at the 5\% level}
\item Reject the null at the 10\% level
\item Fail to reject the null at the 10\% level
\end{opts}

\item Consider the null hypothesis that $H_0:\beta_{educ}=-0.15$. What is the p-value associated with this hypothesis test?
\begin{opts}
\item 0.1998
\item 0.1639
\item 0.1882
\item 0.1510
\item \ca{None of the above}
\end{opts}

\item What is the covariance between the estimate of the parameter on \textit{(Intercept)} and the parameter on \textit{west}?
\begin{opts}
\item $-0.03627$
\item $-0.04409$
\item $-0.04026$
\item \ca{$-0.02994$}
\item $-0.02254$
\end{opts}

\item \multipart{9}{15}
The data set \code{catholic} contains data on students that have recently left high school. You can load the data into R using the \code{wooldridge} library using:

\medskip
\code{library(wooldridge)}\\
\code{dat = wooldridge::catholic}
\medskip

The variable of interest is \textit{hsgrad} which indicates if the student graduated high school. We would like to study what are the main determinants of high school graduation using a linear probability model. To answer this question you will need to add a new variable to the data set called \textit{nonwhite} which should equal 1 if the student is non-white (i.e., asian, hispanic, or black) and zero otherwise.\\
What percent of the population are \textit{nonwhite}?
\begin{opts}
\item 16.48\%
\item \ca{22.58\%}
\item 16.18\%
\item 15.43\%
\item None of the above
\end{opts}

\item Estimate a linear probability model on high school graduation \textit{hsgrad}.\\
Include in your model controls for if the student attended a Catholic high school (\textit{cathhs}), \textit{nonwhite}, and both the 12th grade reading and math scores. Your model should have 5 parameters including the intercept.\\
What is the R-squared from the estimated model?
\begin{opts}
\item 0.03834
\item 0.04439
\item 0.04080
\item \ca{0.07493}
\item 0.05675
\end{opts}

\item What is the correct interpretation of the parameter on \textit{cathhs}? Holding fixed the variables in the model, attending a Catholic high school increases the probability of high school graduation by \dots
\begin{opts}
\item 3.0198 percent
\item 0.3020 percent
\item \ca{3.0198 percentage points}
\item 0.0302 percent
\item None of the above
\end{opts}

\item What is the estimate of the lower bound of the 90\% confidence interval for the intercept of the model?
\begin{opts}
\item 0.3342
\item 0.2780
\item 0.4815
\item \ca{0.5120}
\item None of the above
\end{opts}

\item Consider the null hypothesis that $H_0:\beta_{math12}=0$. What is the strongest conclusion we can draw about this null hypothesis?
\begin{opts}
\item \ca{Reject the null at the 0.1\% level}
\item Reject the null at the 1\% level
\item Reject the null at the 5\% level
\item Reject the null at the 10\% level
\item Fail to reject the null at the 10\% level
\end{opts}

\item Re-estimate the model adding log family income (\textit{lfaminc}) as an additional control. What is the percentage change in the parameter on \textit{cathhs} when we control for family income?
\begin{opts}
\item $-53.62\%$
\item $-51.42\%$
\item 47.80\%
\item $-54.01\%$
\item \ca{None of the above}
\end{opts}

\item Re-estimate the model adding log family income (\textit{lfaminc}) as an additional control. What is the variance inflation factor for the parameter on \textit{cathhs} when we control for family income? (note: if you are calculating the variance inflation factor from the formula you need to ensure that the same observations are used in the different regressions)
\begin{opts}
\item \ca{1.0249}
\item 1.2479
\item 1.2946
\item 0.9104
\item None of the above
\end{opts}
\end{enumerate}

\keyheader
\begin{keytable}
\krow{1}{a}{2.7431}{\code{mean(dat\$kids)}}
\krow{2}{a}{0}{\code{sum(is.na(dat\$age))}}
\krow{3}{e}{$-0.12260$}{\code{lm(kids \textasciitilde\ educ + age + I(age\^{}2) + year + east + northcen + west)}}
\krow{4}{b}{2809.876}{SSR = \code{sum(resid(m)\^{}2)}}
\krow{5}{c}{2.2444}{\code{qt(1-0.025/2, df.residual(m))}, df $=1129-8=1121$}
\krow{6}{c}{$p=0.0237$}{$0.01<p<0.05$ $\Rightarrow$ reject at 5\%}
\krow{7}{e}{0.1351}{$t=(\hat\beta+0.15)/se$, two-sided \code{2*pt(-|t|, df)}}
\krow{8}{d}{$-0.02994$}{\code{vcov(m)["(Intercept)","west"]}}
\krow{9}{b}{22.58\%}{\code{nonwhite = asian|hispan|black}, then \code{mean}}
\krow{10}{d}{0.07493}{$n=5970$; \textit{hsgrad} has 1460 NAs (lm drops them)}
\krow{11}{c}{0.0302}{LPM coefficient $\times100$ = percentage points}
\krow{12}{d}{0.5120}{\code{confint(m1, level=0.90)}}
\krow{13}{a}{$p\approx10^{-43}$}{$p<0.001$}
\krow{14}{e}{$-39.46\%$}{\code{100*(b\_new/b\_old - 1)}}
\krow{15}{a}{1.0249}{\code{car::vif(m2)}; by hand only on rows with non-NA \textit{hsgrad} (full sample gives 1.0167)}
\end{keytable}

\subsection*{R code}
\lstinputlisting{../key_original.r}
\end{document}
